\section{Formalization}\label{sec:lean}

\paragraph{What is checked.} The general statements of all thirteen entries are checked in Lean~4
\citep{DeMouraUllrich2021}, in a library with no external dependencies, together with the Collatz
ledger and split pairs (\Cref{thm:affine-ledger,thm:split-pair}), the four-phase family and the
orbit of $22$ (\Cref{thm:four-phase,cor:orbit-22}), Ducci nilpotency (\Cref{thm:ducci}), the two
separating examples for G8, and the instances in \Cref{ex:G4-intervals,ex:G10-checksum,ex:G13-powers}
and \Cref{prop:rule184}, the binary counter of \Cref{ex:G3-counter}, a three-site sandpile, the
colouring obstruction of the Moser spindle and the integer Ducci bound for $n = 4$. The last column of
\Cref{tab:catalog} says, entry by entry, what is covered; sandpiles on arbitrary graphs, the
plane embedding of the Moser spindle, the other remarks and the imported results are proved on paper
only. Some Lean statements are formally weaker than the printed ones, for
example \Cref{thm:G2} without the extraction of a terminal stage, and some carry a hypothesis they
do not use; both are listed in \Cref{supp:status}.

\paragraph{Axioms.} The library builds without \texttt{sorry} or added axioms and uses at most the
standard axioms \texttt{propext}, \texttt{Quot.sound} and \texttt{Classical.choice}.

\paragraph{Where to look.} Following \citet{Pollack1998}, what a reader must trust is the kernel and
the statements. \Cref{supp:concordance} maps every numbered result of this paper to its Lean
declarations, file and line; \Cref{supp:status} gives the axioms of each declaration and a
clause-by-clause comparison with the printed statements; \Cref{supp:replay} gives the commands that
rebuild the library.
